\section{Các nghiên cứu liên quan}

Các nghiên cứu liên quan chỉ tập trung vào ba câu hỏi trực tiếp dẫn đến đề tài: FoodSeg103 khó ở đâu, các mô hình segmentation nhẹ đang đánh đổi điều gì, và các cơ chế tăng cường chi tiết cục bộ có thật sự gợi mở một hướng phù hợp cho dữ liệu thực phẩm hay không.

\subsection{Fine-grained food segmentation và benchmark FoodSeg103}

FoodSeg103 là benchmark ingredient-level tiêu biểu cho bài toán này \cite{benchmarkfood}. Bộ dữ liệu gồm 7{,}118 ảnh, 103 lớp nguyên liệu và 42{,}097 mask; trong đó tập huấn luyện có 4{,}983 ảnh và tập đánh giá có 2{,}135 ảnh. Kết quả benchmark ban đầu cho thấy DeepLabv3+ chỉ đạt 34.2 mIoU trên FoodSeg103, thấp hơn đáng kể so với khi đánh giá trên các bộ dữ liệu food segmentation đơn giản hơn \cite{benchmarkfood}. Điều đó cho thấy khó khăn chính không nằm ở việc nhận diện món ăn nói chung, mà ở khả năng tách chính xác từng nguyên liệu có tín hiệu phân biệt rất mịn.

Nhóm công trình này làm tốt ở chỗ xác lập một benchmark thực sự khó và chỉ ra vai trò của semantic prior, ví dụ ReLeM giúp SeTR tăng từ 41.3 lên 43.9 mIoU \cite{benchmarkfood}. Tuy nhiên, các hướng mạnh về biểu diễn hoặc tri thức đa phương thức vẫn chưa giải quyết trực tiếp bài toán triển khai nhẹ. Vì vậy, giá trị lớn nhất của nhánh này đối với đề tài hiện tại là làm rõ rằng FoodSeg103 đòi hỏi mô hình phải nhạy với texture, boundary và class imbalance, chứ không chỉ mạnh về ngữ cảnh toàn cục.

\subsection{Lightweight semantic segmentation cho triển khai hiệu quả}

Nhóm công trình thứ hai tập trung vào trade-off giữa độ chính xác và chi phí suy luận. Các mô hình như BiSeNet, Fast-SCNN và PIDNet sử dụng backbone nhẹ, downsampling chiến lược hoặc kiến trúc đa nhánh để giảm tham số, FLOPs và độ trễ \cite{bisenet,poudel2019fastscnn,xu2023pidnet}. Trên Cityscapes, BiSeNet đạt 68.4\% mIoU ở 105 FPS, Fast-SCNN đạt 68.0\% mIoU ở 123.5 FPS với 1.11M tham số, còn PIDNet-S đạt 78.6\% mIoU ở 93.2 FPS với 7.6M tham số \cite{bisenet,poudel2019fastscnn,xu2023pidnet}.

Điểm mạnh của nhóm này là chứng minh rằng segmentation thời gian thực hoàn toàn khả thi nếu kiến trúc được tổ chức hợp lý. Đặc biệt, BiSeNet đáng chú ý vì tách riêng nhánh giữ chi tiết không gian và nhánh học ngữ cảnh, nên phù hợp để làm baseline khi mục tiêu là phân tích rõ phần nào của kiến trúc còn thiếu cho dữ liệu fine-grained.

Tuy nhiên, điểm yếu của lightweight segmentation cũng khá rõ. Phần lớn lợi ích về tốc độ đến từ việc nén biểu diễn hoặc giảm độ phân giải đặc trưng sớm, trong khi đó lại là thao tác dễ làm mất texture và boundary cues. Trên các benchmark cảnh đường phố, mất mát này đôi khi còn chấp nhận được vì lớp vẫn có cấu trúc hình học lớn và ngữ cảnh mạnh. Trên FoodSeg103 thì ngược lại: mô hình nhẹ có nguy cơ bỏ mất đúng loại tín hiệu mà bài toán cần nhất. Vì vậy, lightweight models là nền tảng tốt cho triển khai, nhưng tự thân chúng chưa đủ cho ingredient-level food segmentation.

\subsection{Các cơ chế tăng cường chi tiết cục bộ và texture-aware}

Để bù lại phần thông tin cục bộ bị mất trong các kiến trúc gọn nhẹ, nhiều nghiên cứu đề xuất các cơ chế detail-enhanced hoặc texture-aware. Gated-SCNN học thông tin biên qua một shape stream riêng \cite{takikawa2019gatedscnn}; PIDNet kết hợp ba nhánh detail, context và boundary để cải thiện mask mà không phình to backbone quá mức \cite{xu2023pidnet}; PointRend chỉ tinh chỉnh tại các điểm bất định cao quanh ranh giới \cite{kirillov2020pointrend}; còn Deep-TEN cho thấy texture có thể được biểu diễn như một tín hiệu thống kê có cấu trúc thay vì chỉ là pattern cục bộ rời rạc \cite{zhang2017deepten}.

Điểm chung của các hướng này là chúng không cố giải quyết vấn đề bằng cách làm mạng lớn hơn, mà tổ chức lại nơi nào cần giữ chi tiết và nơi nào cần refinement. Đây là một lập luận rất phù hợp với ảnh thực phẩm, nơi sai số quan trọng thường tập trung ở vùng giao nhau giữa các nguyên liệu hoặc ở các lớp có bề mặt gần giống nhau.

Dù vậy, bằng chứng thực nghiệm cho các cơ chế này dưới ràng buộc \emph{lightweight food segmentation} vẫn còn hạn chế. Phần lớn chúng được kiểm chứng trên benchmark tổng quát hoặc đô thị, nơi biên rõ và cấu trúc hình học ổn định hơn. Vì thế, khoảng trống nghiên cứu không phải là thiếu ý tưởng detail-aware, mà là thiếu một pipeline baseline đủ sạch để kiểm tra có hệ thống xem các ý tưởng đó có thật sự cải thiện trade-off accuracy--efficiency trên FoodSeg103 hay không.

\subsection{Khoảng trống nghiên cứu và vị trí của bài hiện tại}

Từ ba nhóm công trình trên có thể rút ra một kết luận chung. Các mô hình mạnh trên FoodSeg103 cho thấy bài toán cần ngữ cảnh và biểu diễn tốt; các mô hình nhẹ cho thấy triển khai hiệu quả là khả thi; còn các hướng detail-aware cho thấy local refinement là một giả thuyết hợp lý để giảm mất mát chi tiết. Nhưng ba hướng này chưa được nối với nhau một cách chặt chẽ trên cùng một thiết lập baseline nhẹ và dữ liệu đã được xử lý ổn định.

Do đó, bài hiện tại không cố tuyên bố một mô-đun texture-aware mới đã vượt trội hơn state of the art. Thay vào đó, nghiên cứu tập trung vào bước nền: phân tích dữ liệu và chuẩn hóa baseline BiSeNet trên \texttt{FoodSeg103} nguyên thủy. Song song, một nhánh thực nghiệm mở rộng về tái cấu trúc không gian nhãn được dùng để phân tích bổ sung ảnh hưởng của phân mảnh nhãn. Thiết kế hai tầng này tạo điểm tựa cần thiết để các nghiên cứu tiếp theo đánh giá công bằng các cải tiến texture-aware hoặc boundary-aware dưới cùng điều kiện thực nghiệm.
